\subsection*{电荷、叠加原理与连续分布}

普通物质由原子组成，原子核带正电、电子带负电。宏观带电通常是电子转移的结果，而不是产生了
新的电荷。电荷具有代数可加性并满足守恒；孤立系统总电荷不随内部过程改变。微观上电荷以元电荷
$e$ 为基本单位量子化，质子带 $+e$、电子带 $-e$。宏观物体所含带电粒子数极多，才常把电荷分布
近似为连续函数。

真空中两个静止点电荷的相互作用满足
\[
 \vec F_{12}=\frac1{4\pi\varepsilon_0}
 \frac{q_1q_2}{r_{12}^3}\vec r_{12}.
\]
矢量形式同时包含大小和方向，同号相斥、异号相吸。多个电荷共同作用时，先分别求每个电荷产生
的力或电场，再作矢量和，这就是线性叠加原理。连续分布则把求和变为积分：
\[
 \vec E(\vec r)=\frac1{4\pi\varepsilon_0}
 \int\frac{\vec r-\vec r'}{|\vec r-\vec r'|^3}\,\dd q'.
\]
选择线、面或体电荷元并不是形式问题，而由带电体几何与观察尺度决定。

\subsection*{场线、电通量与高斯定律}

电场定义为单位正试探电荷所受力，$\vec E=\vec F/q$。试探电荷应足够小，避免显著改变原来的
电荷分布。电场线在每点的切线方向给出 $\vec E$ 方向，疏密只定性表示强弱；静电场线从正电荷
或无穷远出发，终止于负电荷或无穷远，彼此不相交，也不形成闭合曲线。场线是一种图示工具，
并非真实物质轨迹。

通过有向面元的电通量为
\[
 \dd\Phi_E=\vec E\cdot\dd\vec S,
 \qquad
 \Phi_E=\oint_S\vec E\cdot\dd\vec S.
\]
闭合面的面元规定向外。高斯定律
\[
 \oint_S\vec E\cdot\dd\vec S=\frac{Q_{\rm in}}{\varepsilon_0}
\]
对任意闭合面成立，但只有电荷分布具有球、柱或平面对称时，才能把场强从积分号中直接提出。
面外电荷虽不贡献净通量，却仍可能在面上产生非零电场；“净通量为零”也不等于“处处无场”。

球对称分布外部的场只由总电荷决定，等效为点电荷；均匀带电实心球内部包围电荷随 $r^3$ 增长，
故场强随 $r$ 线性增长。无限长均匀线电荷有
$E=\lambda/(2\pi\varepsilon_0r)$，无限大均匀薄板单侧有
$E=\sigma/(2\varepsilon_0)$。有限物体只能在离边缘足够远或观察距离满足相应条件时近似使用这些
结果，不能把理想对称结论无条件搬到边缘附近。

偶极子由相距很近的异号等量电荷构成，偶极矩 $\vec p=q\vec d$ 从负电荷指向正电荷。远处电场
按 $1/r^3$ 衰减，比点电荷的 $1/r^2$ 更快；总电荷为零并不意味着场处处为零，只是最低阶的单极项
消失。更复杂电荷分布可依次用单极、偶极、四极等多极矩描述其远场。

